\documentclass[12pt]{article}
\usepackage{amsmath,amssymb}
\usepackage[margin=1in]{geometry}
\usepackage{graphicx}

\title{Velocity and Acceleration Analysis}
\author{}
\date{}

\begin{document}

\maketitle

\section*{Problem Description}

This document presents the velocity and acceleration analysis of a four-bar linkage mechanism using both graphical (Asymptote diagrams) and analytical methods.

\section*{Generated Diagrams}

The analysis produces three separate diagrams from a single Asymptote file (\texttt{q3-2.asy}):

\begin{enumerate}
  \item \textbf{q3-2-space.pdf} — Mechanism geometry showing the spatial configuration of all links and pivot points
  \item \textbf{q3-2-velocity.pdf} — Velocity polygon diagram using graphical analysis methods
  \item \textbf{q3-2-acceleration.pdf} — Acceleration polygon diagram using graphical analysis methods
\end{enumerate}

\section*{Analysis Results}

\subsection*{Point B Location}

The slider point $B$ is located at:
$$B = (228.7969195251, 661.5526960234) \text{ mm}$$

\subsection*{Velocity Analysis}

The input link (link 2) rotates at a constant angular velocity (60 RPM).

\begin{table}[h]
\centering
\begin{tabular}{|l|r|}
\hline
\textbf{Parameter} & \textbf{Value} \\\\
\hline
$\omega_2$ (Input angular velocity) & 6.28318530717959 rad/s \\\\
$v_{A_2/O_2}$ (Velocity at $A$ on link 2) & 1099.5574287564275 mm/s \\\\
$v_{A_4/O_4}$ (Velocity at $A$ on link 4) & 910.73802522663527 mm/s \\\\
$v_{A_2/A_4}$ (Relative velocity) & 616.10290418057082 mm/s \\\\
$\omega_4$ (Angular velocity of link 4) & 6.57220835691986 rad/s \\\\
\hline
\end{tabular}
\caption{Velocity Analysis Results}
\end{table}

\subsection*{Acceleration Analysis}

\begin{table}[h]
\centering
\begin{tabular}{|l|r|}
\hline
\textbf{Parameter} & \textbf{Value} \\\\
\hline
$a^n_{A_2/O_2}$ (Centripetal acceleration at $A$ on link 2) & 6908.72308076254831 mm/s$^2$ \\\\
$a^n_{A_4/O_4}$ (Centripetal acceleration at $A$ on link 4) & 5985.56006035918017 mm/s$^2$ \\\\
$a^c_{A_2/A_4}$ (Coriolis acceleration) & 8098.31331115628382 mm/s$^2$ \\\\
$a^t_{A_4/O_4}$ (Tangential acceleration at $A$ on link 4) & 11969.40202641432734 mm/s$^2$ \\\\
$\alpha_4$ (Angular acceleration of link 4) & 86.37544699614139 rad/s$^2$ \\\\
$a_{A_4/O_4}$ (Total acceleration at $A$ on link 4) & 13382.58249016602895 mm/s$^2$ \\\\
\hline
\end{tabular}
\caption{Acceleration Analysis Results}
\end{table}

\subsection*{Slider Point B Acceleration}

At the slider point $B$:

\begin{table}[h]
\centering
\begin{tabular}{|l|r|}
\hline
\textbf{Component} & \textbf{Value} \\\\
\hline
$a^n_B$ (Normal acceleration component) & 30235.74588073705308 mm/s$^2$ \\\\
$a^t_B$ (Tangential acceleration component) & 60462.81289729897253 mm/s$^2$ \\\\
$a_B$ (Total acceleration magnitude) & 67601.42063905381656 mm/s$^2$ \\\\
\hline
\end{tabular}
\caption{Point B Acceleration}
\end{table}

\section*{Key Relationships}

The analysis is based on the following kinematic relationships:

\textbf{Velocity:}
\begin{align*}
v &= \omega \cdot r \\\\
v_{rel} &= v_1 - v_2
\end{align*}

\textbf{Acceleration (components):}
\begin{align*}
a^n &= \omega^2 \cdot r \quad \text{(centripetal)} \\\\
a^t &= \alpha \cdot r \quad \text{(tangential)} \\\\
a^c &= 2\omega_1 \cdot v_{rel} \quad \text{(Coriolis)} \\\\
a_{total} &= \sqrt{(a^n)^2 + (a^t)^2}
\end{align*}

\section*{Notes}

\begin{itemize}
  \item All distances are measured in millimeters (mm)
  \item All velocities are in mm/s
  \item All accelerations are in mm/s$^2$
  \item All angular quantities are in radians per second (rad/s) or rad/s$^2$
  \item The mechanism analysis assumes constant input angular velocity
  \item Point $A$ is the connection between links 2 and 4
  \item Point $B$ is the slider constrained to move along link 4
\end{itemize}

\end{document}